\documentclass[10pt,a4paper]{article}
\usepackage[utf8]{inputenc}
\usepackage[T1]{fontenc}
\usepackage[margin=0.6in]{geometry}
\usepackage{amsmath,amssymb}
\usepackage{hyperref}

% Professional Links Color
\usepackage{xcolor}
\definecolor{darkblue}{RGB}{2, 132, 199}
\hypersetup{colorlinks=true, linkcolor=darkblue, urlcolor=darkblue}

% Disable paragraph indentation and page numbers for a clean look
\setlength{\parindent}{0pt}
\pagestyle{empty}

% Easy-to-maintain custom section command
\newcommand{\cvsection}[1]{%
    \vspace{14pt}
    {\large\bfseries\MakeUppercase{#1}}
    \vspace{2pt}\hrule\vspace{6pt}
}

\begin{document}

% ==========================================
% HEADER
% ==========================================
\begin{center}
    {\LARGE \bfseries Joan Alcaide-Núñez} \\[4pt]
    \small
    Email: \href{mailto:joan.alcaide@campus.lmu.de}{joan.alcaide@campus.lmu.de} $\cdot$
    Website: \href{https://joanalnu.github.io}{joanalnu.github.io} $\cdot$
    GitHub: \href{https://github.com/joanalnu}{github.com/joanalnu}
\end{center}

% ==========================================
% EDUCATION
% ==========================================
\cvsection{Education}

\textbf{Ludwig-Maximilians-Universität (LMU) Munich} \hfill Munich, Germany \\
B.Sc. in Physics \hfill Oct 2025 -- Present \\
\textit{Core concentration in foundational and mathematical physics pathways.}

\vspace{6pt}

\textbf{German School of Barcelona} \hfill Barcelona, Spain \\
Dual High School Diploma (Germany + Spain) \hfill Sept 2017 -- May 2025 \\
\textit{Academic Achievement Award with a perfect Grade Point Average of 1.0 / 1.0.}

\vspace{6pt}

\textbf{Additional Specialised Training:}
\begin{itemize}
    \itemsep0em \topsep0pt \parsep2pt \partopsep0pt
    \item \textit{Gravity and Black Holes} -- Perimeter Institute for Theoretical Physics (GoPhysics!), 2025
    \item \textit{Mathematics, Statistics, and Data Science} -- Autonomous University of Barcelona (UAB), 2022
    \item \textit{Algorithms and Programming (C++)} -- Polytechnic University of Catalonia (UPC), 2022
    \item \textit{League of Codes (C++ Masterclass)} -- Harbour Space University, 2022 -- 2023
\end{itemize}

% ==========================================
% RESEARCH EXPERIENCE
% ==========================================
\cvsection{Research Experience}

\textbf{GRB Cosmography \& Relativistic Foundations} \hfill Merate, Italy \\
Osservatorio Astronomico di Brera (OAB-INAF) \hfill June -- July 2025
\begin{itemize}
    \itemsep0em \topsep2pt
    \item Constrained global cosmological parameters ($\Omega_m, \Omega_\Lambda$) via $\chi^2$ minimisation profiles using the empirical $E_{\text{peak}}-E_{\text{iso}}$ relation of Gamma-Ray Bursts.
    \item Modeled detection synergies with mock binary neutron star populations matching Einstein Telescope specifications.
    \item Received foundational technical training on relativistic jet kinematics from Dr. Gabriele Ghisellini.
    \item \textit{Supervisors:} Dr. Giancarlo Ghirlanda, Dr. Om Sharan Salafia $\cdot$ \href{https://github.com/joanalnu/oab-inaf}{Repository Link}.
\end{itemize}

\textbf{CAPIBARA Project} \hfill International \\
Founder \& Project Lead -- CAPIBARA Collaboration \hfill 2024 -- 2026
\begin{itemize}
    \itemsep0em \topsep2pt
    \item Established the CAPIBARA Collaboration to connect science-driven international students across secondary and university levels to coordinate collaborative learning frameworks.
    \item Leading the \href{https://capibara3.github.io/gamma/}{GAMMA Team} evaluating the technical feasibility, instrumentation parameters, and mission profiles for a student-designed gamma-ray transient observing CubeSat.
\end{itemize}

\textbf{Type Ia Supernovae: Mapping Expansion History} \hfill Barcelona, Spain \\
Institute of Space Sciences (ICE-CSIC-IEEC) \hfill June -- July 2024
\begin{itemize}
    \itemsep0em \topsep2pt
    \item Extracted infrared imaging data from ESO (SOFI) and cross-matched optical photometry from ZTF and ATLAS using the \texttt{sncosmo} Python suite to map luminosity distances.
    \item Explored non-linear expansion, solving the Friedmann acceleration equations to fit cosmic density parameters via light-curve parameterisation.
    \item \textit{Supervisor:} Dr. Lluís Galbany.
\end{itemize}

\textbf{Cepheid Variable Stars Research} \hfill Budapest, Hungary \\
Research Fellow -- Youth \& Science Program \hfill Aug -- Dec 2023
\begin{itemize}
    \itemsep0em \topsep2pt
    \item Reconstructed classical Hubble-Lemaître expansion constraints, deriving an empirical value of $H_0 = 70.26 \pm 7.1 \, \text{km} \, \text{s}^{-1} \text{Mpc}^{-1}$ using modern Cepheid observations from the Konkoly Observatory.
    \item Selected for a selective 3-year funded research track, marking first formal exposure to academic data review $\cdot$ \href{https://joanalnu.github.io/files/Paper2023_RecomputingTheHubbleConstant.pdf}{Technical Report}.
    \item \textit{Supervisors:} Dr. Ignasi Pérez-Ràfols, Dr. Laia Casamiquela.
\end{itemize}

\textbf{German Aerospace Center (DLR)} \hfill Oberpfaffenhofen, Germany \\
Research Intern -- Institute for Remote Sensing (DLR-IMF) \hfill Oct 2024
\begin{itemize}
    \itemsep0em \topsep2pt
    \item 2-week internship at the photogrammetry (hyperspectral imaging) and satellite observations (thermal LandSat-8 imaging python-based analysis) teams
\end{itemize}

% ==========================================
% AWARDS & HONORS
% ==========================================
\cvsection{Awards \& Honors}

\textbf{Youth \& Science Fellowship} $\cdot$ \textit{Fundació Catalunya La Pedrera} \hfill 2023 -- 2026 \\
Fully funded 3-year international research fellowship supporting consecutive competitive summer research placements.

\vspace{4pt}

\textbf{Abiturpreis (Physics Distinction)} $\cdot$ \textit{German Physical Society (DPG)} \hfill 2025 \\
National prize awarded for outstanding academic excellence and foundational focus in physics during high school.

\vspace{4pt}

\textbf{International Astronomy and Astrophysics Competition (IAAC)} \hfill 2023 -- 2025 \\
Awarded the \textbf{Silver Honour} and the \textbf{National Award for Spain} (2025), ranking in the top 7\% globally out of more than 12,000 international participants.

\vspace{4pt}

\textbf{Jugend Forscht (Youth Research Competition)} \hfill 2024 -- 2025 \\
Awarded \textbf{First Prize Iberia} across regional and national tiers; received the Special Award in Scientific Photography (NRW, 2025).

\textbf{High School Competitions} \hfill

participation in a lsit of of mathematics, coding and physics competitions

% ==========================================
% OUTREACH & TALKS
% ==========================================
\cvsection{Talks \& Science Outreach}

\begin{itemize}
    \itemsep0em \topsep0pt
    \item \textbf{SSEA Symposium (April 2026):} Presented \textit{``The CRD Instrument and the COSMOS Program.''}
    \item \textbf{CosmoXarxa (March 2025):} Delivered an invited presentation on high-energy astrophysics metrics at the Barcelona Science Museum.
    \item \textbf{Academic Workshops (2023 -- 2024):} Designed introductory astrophysics workshops for 11th graders at the German School Barcelona and organized planetary space exploration seminars for elementary tracks.
    \item \textbf{Museum Explainer (2022 -- 2023):} Content volunteer conducting interactive public physics experiments at the CosmoCaixa Science Museum.
\end{itemize}

% ==========================================
% WORK EXPERIENCE
% ==========================================
\cvsection{Work Experience}

\textbf{Ludwig-Maximilians-Universität (LMU) Munich} \hfill Munich, Germany \\
Student Assistant -- Central Diversity Management \& Research Funding Unit \hfill Oct 2025 -- Present

% ==========================================
% TECHNICAL SKILLS
% ==========================================
\cvsection{Technical Skills}

\textbf{Programming \& Core Frameworks:} Python (\textit{NumPy, SciPy, Pandas, Matplotlib, Astropy}), \texttt{sncosmo}, C++, \LaTeX, Markdown. \\
\textbf{Tools \& Environments:} Git, GitHub, Visual Studio Code, NeoVim, Unix/Linux command-line environments.

\end{document}
